\scanpage{32}%
\mspage{17}%
et la loi de composition par
\[
(u',s',t')\mathbin{\handwritten{o}}(u,s,t)
=(u\,u',u's+t',u^{\raisebox{-.38ex}{$\scriptstyle\prime$}\,3}t+t'+u^{\raisebox{-.38ex}{$\scriptstyle\prime$}\,2}rs')\, .
\]

Lorsque la courbe est mise sous forme (6.8), ces équations se réduisent à
\[
\begin{aligned}
u\in\mu_3 \qquad s\in\mathbb{F}_4 \qquad t^2+t=s^3
 &\quad ({}=0\ \text{ou}\ 1\ \text{selon que}\\
 &\qquad s=0\ \text{ou}\ s\ne 0) \, .
\end{aligned}
\]

7. \quad \uline{Anneau des formes modulaires entières}.

On appelle forme modulaire entière de poids $n$ une loi qui à chaque courbe de genre un sur une base $S$ associe une section de $\omega^{\otimes n}$, et ce de façon compatible au changement de base. Il reviendrait\margin{\handwritten{\tikz[baseline=-0.2ex]{\draw[thin,gray] (0,-0.35em) -- (0.48em,1.30em);} \uuline{pourquoi} ?}} au même de se limiter aux courbes sans fibres additives.

\uline{Proposition} 7.1. \quad \uline{L'anneau des formes modulaires entières est engendré sur} $\mathbb{Z}$ \uline{par} $c_4$, $c_6$ \uline{et} $\Delta$ \uline{soumis à la seule relation}
\[
c_4^3-c_6^2=1728\,\Delta \, . \tag*{(7.2.)}
\]
\margin{\tikz[baseline=(note.north)]{\node[inner sep=2pt,text width=9.8em,align=left] (note) {\handwritten{N.B. Sur le site modulaire, les diviseurs de $c_4$, $c_6$, $\Delta$ sont à \unsure{multiplicité|multiplicités} $1$ ; Celui de $\Delta$ est lisse sur $\mathbb{Z}$, ceux de $c_4$, $c_6$ le sont sauf au dessus des points $2$ et $3$.}};\handwritten{\draw[thin,rounded corners=1pt] (note.north west) -- (note.south west) -- (note.south east);}}}

Appliquant la définition à la courbe (1.3), on voit qu'une forme modulaire entière de poids $n$ s'exprime comme poly\tikz[baseline=(nome.base)]{\node[inner sep=0pt] (nome) {nôme};\handwritten{\draw[thin,gray,rotate around={-25:(nome.center)}] (nome.center) ellipse (1.04em and 1.65em);}} de poids $n$ en les $a_i$. Pour $2$ et $3$ inversibles, il résulte de 2.5 que ces formes sont \tikz[baseline=(des.base)]{\node[inner sep=0pt] (des) {des};\handwritten{\draw[thin,gray,rotate around={-15:(des.center)}] (des.center) ellipse (1em and 1.02em);}} poly\tikz[baseline=(noms.base)]{\node[inner sep=0pt] (noms) {nôm};\handwritten{\draw[thin,gray,rotate around={-25:(noms.center)}] (noms.center) ellipse (0.84em and 1.47em);}}es en $c_4$ et $c_6$, et en particulier que $c_4$, $c_6$ et $\Delta$ ne satisfont pas d'autres relations que (7.2). Reste à prouver que si un poly\tikz[baseline=(nomt.base)]{\node[inner sep=0pt] (nomt) {nôme};\handwritten{\draw[thin,gray,rotate around={-25:(nomt.center)}] (nomt.center) ellipse (0.85em and 1.53em);}} en $c_4$, $c_6$ et $\Delta$ est identiquement nul en caractéristique $2$ (resp. $3$), alors il est divisible par $2$ (resp.$3$) dans $\mathbb{Z}[c_4,c_6,\Delta]/\handwritten{\bigl(}c_4^3-c_6^2-1728\,\Delta\handwritten{\bigr)}$, ce qui signifie%
% archive sheet number: 31
